\section{Normalization Validation}

\subsection{Functional Dependencies and Candidate Keys}

The following dependencies are derived from the implemented identifiers and uniqueness constraints. The phrase \emph{all non-key attributes} means every remaining descriptive attribute of that relation.

\small
\begin{longtable}{>{\raggedright\arraybackslash}p{3.6cm}>{\raggedright\arraybackslash}p{4.1cm}>{\raggedright\arraybackslash}p{6.8cm}}
  \caption{Candidate keys and principal functional dependencies.}
  \label{tab:functional-dependencies}\\
  \toprule
  \textbf{Relation} & \textbf{Candidate key(s)} & \textbf{Non-trivial functional dependencies}\\
  \midrule
  \endfirsthead
  \toprule
  \textbf{Relation} & \textbf{Candidate key(s)} & \textbf{Non-trivial functional dependencies}\\
  \midrule
  \endhead
  \nolinkurl{ROLE} &
  role\_id; role\_name &
  role\_id $\rightarrow$ role\_name; role\_name $\rightarrow$ role\_id\\
  \nolinkurl{USERS} &
  user\_id; email &
  user\_id $\rightarrow$ all non-key attributes; email $\rightarrow$ all other attributes\\
  \nolinkurl{SPACE_TYPE} &
  space\_type\_id; space\_type\_name &
  space\_type\_id $\rightarrow$ space\_type\_name; space\_type\_name $\rightarrow$ space\_type\_id\\
  \nolinkurl{SPACES} &
  space\_code; (building, room\_number) &
  Each candidate key $\rightarrow$ all non-key attributes\\
  \nolinkurl{AUTO_USAGE_POLICY} &
  (space\_type\_id, role\_id) &
  No non-key attributes; only trivial dependencies\\
  \nolinkurl{FACILITY} &
  facility\_id &
  facility\_id $\rightarrow$ space\_code, facility\_name, description\\
  \nolinkurl{BOOKING_REQUEST} &
  booking\_id &
  booking\_id $\rightarrow$ all non-key attributes\\
  \nolinkurl{BOOKING_DECISION} &
  decision\_id; booking\_id &
  decision\_id $\rightarrow$ all non-key attributes; booking\_id $\rightarrow$ all other attributes\\
  \nolinkurl{USAGE_SESSION} &
  session\_id; decision\_id &
  session\_id $\rightarrow$ all non-key attributes; decision\_id $\rightarrow$ all other attributes\\
  \nolinkurl{MAINTENANCE_RECORD} &
  maintenance\_id &
  maintenance\_id $\rightarrow$ all non-key attributes\\
  \nolinkurl{BOOKING_NOTIFICATION} &
  (booking\_id, maintenance\_id, notification\_type) &
  Composite key $\rightarrow$ notification\_time\\
  \bottomrule
\end{longtable}
\normalsize

Foreign keys do not create reverse functional dependencies. For example, a role can belong to many users and a space can have many requests. Conditional rules such as \emph{automatic decisions have no deciding staff member} are business constraints, not additional functional dependencies.

\subsection{Normal-Form Proof}

\textbf{First Normal Form.} All implemented columns contain atomic scalar values. Repeating sets such as required facilities and automatic policies are represented by separate rows and relations.

\textbf{Second Normal Form.} Relations with single-column keys have no possible partial dependency. \texttt{AUTO\_USAGE\_POLICY} has no non-key attributes. In \texttt{BOOKING\_NOTIFICATION}, \texttt{notification\_time} describes one complete booking--maintenance--notification-type event and depends on the full composite key.

\textbf{Third Normal Form.} Every non-trivial dependency listed in Table~\ref{tab:functional-dependencies} has a candidate key as its determinant. No non-key attribute determines another non-key attribute within the same relation. Lookup descriptions such as role name and space-type name are stored only in their lookup relations, eliminating the Phase 1 transitive dependencies.

Consequently, every relation satisfies 3NF. Under the documented dependency set, every non-trivial determinant is also a superkey, so the relations satisfy the stronger Boyce--Codd Normal Form as well.

\section{LLM and Agent Improvement Process}

\subsection{Model and Framework}

The project used \textbf{DeepSeek V4 Flash} through the \textbf{OpenCode} framework. The LLM was used as a research and drafting assistant within a human-supervised process; final deliverables were reviewed and finalized by the group.

The framework separates reusable knowledge from experiment artifacts:

\begin{itemize}
  \item the Global Agent stores project-wide terminology, invariants, consistency rules, and verification behavior;
  \item section-specific Skills store task methods and checklists;
  \item experiment results provide material for evaluation rather than becoming submissions directly; and
  \item human evaluation determines which reusable improvements are accepted.
\end{itemize}

\subsection{Improvement Cycle}

\begin{enumerate}
  \item Read the specification, current Agent knowledge, and relevant Skill.
  \item Generate an experiment result for one deliverable.
  \item Evaluate requirement coverage, naming consistency, SQL correctness, and unsupported assumptions.
  \item Record the weakness, root cause, and reusable correction.
  \item Refine the relevant knowledge and regenerate.
  \item Human-review the final repository artifact.
\end{enumerate}

For Phase 2, the review emphasis expanded in three directions. Concurrency claims required both an unsafe reproduction and a protected two-session result. Performance claims required repeatable measurements and execution-plan evidence rather than relying only on optimizer suggestions. Requirement traceability was extended from the ERD through migration, maintenance notification, approval locking, and the four required analytical queries.

\begin{table}[H]
  \centering
  \small
  \begin{tabularx}{\textwidth}{>{\raggedright\arraybackslash}p{3.8cm}X}
    \toprule
    \textbf{Evaluation concern} & \textbf{Reusable improvement}\\
    \midrule
    Concurrent correctness &
    State the invariant, reproduce the race, define one lock order, and verify the final database state.\\
    Migration safety &
    Populate replacement columns and tables before removing Phase 1 structures; verify preserved counts.\\
    Large-data validity &
    Make generation deterministic, namespace generated identifiers, and assert coverage/invariants inside the script.\\
    Query performance &
    Compare identical semantics before and after indexing using repeated reads, timings, and plan operators.\\
    Index discipline &
    Evaluate selectivity, storage, and write cost before accepting a missing-index recommendation.\\
    \bottomrule
  \end{tabularx}
  \caption{Representative Phase 2 agent-improvement lessons.}
\end{table}
